% Part: first-order-logic
% Chapter: model-theory
% Section: theory-of-m

\documentclass[../../../include/open-logic-section]{subfiles}

\begin{document}

\olfileid{mod}{bas}{thm}

\section{De theorie van een \printtoken{S}{structure}}

In elke !!{structure}~$\Struct{M}$ zijn sommige !!{sentence}s waar en
andere onwaar. De verzameling van alle !!{sentence}s die zij waar maakt,
heet haar \emph{theorie}. Die verzameling is inderdaad een theorie: elk
gevolg ervan moet in al haar modellen waar zijn, dus ook in~$\Struct{M}$.

\begin{defn}
  Gegeven !!a{structure}~$\Struct M$ is de \emph{theorie} van
  $\Struct{M}$ de verzameling $\Theory{M}$ van !!{sentence}s
  die in $\Struct{M}$ waar zijn, dus $\Theory{M} =
  \Setabs{!A}{\Sat{M}{!A}}$.
\end{defn}

We gebruiken de term ``theorie'' ook informeel voor verzamelingen
!!{sentence}s met een bedoelde interpretatie, ongeacht of zij deductief
gesloten zijn.

\begin{prop}
Voor elke $\Struct{M}$ is $\Theory{M}$ compleet.
\end{prop}

\begin{proof}
Voor elke !!{sentence}~$!A$ geldt ofwel $\Sat{M}{!A}$, ofwel
$\Sat{M}{\lnot !A}$. Dus geldt ofwel $!A \in \Theory{M}$, ofwel
$\lnot !A \in \Theory{M}$.
\end{proof}

\begin{prop}\ollabel{prop:equiv}
  Als $\Struct{N} \models !A$ voor elke $!A \in \Theory{M}$, dan
  $\Struct{M} \elemequiv \Struct{N}$.
\end{prop}

\begin{proof}
Omdat $\Sat{N}{!A}$ voor alle $!A \in \Theory{M}$, geldt
$\Theory{M} \subseteq \Theory{N}$. Als $\Sat{N}{!A}$, dan
$\Sat/{N}{\lnot !A}$, dus $\lnot !A \notin \Theory{M}$. Omdat
$\Theory{M}$ compleet is, geldt $!A \in \Theory{M}$. Bijgevolg
$\Theory{N} \subseteq \Theory{M}$ en dus
$\Struct{M} \elemequiv \Struct{N}$.
\end{proof}

\begin{rem}\ollabel{remark:R}
  Beschouw $\Struct{R} = \langle\Real, <\rangle$: de !!{structure} met
  als domein de verzameling $\Real$ van de reële getallen, in de
  !!{language} die slechts één 2-plaatsig !!{predicate} bevat, geïnterpreteerd
  als de relatie $<$ op de reële getallen. De structuur $\Struct{R}$ is
  duidelijk !!{nonenumerable}. Omdat $\Theory{R}$ echter consistent is,
  heeft zij volgens de stelling van L\"owenheim--Skolem !!a{enumerable}
  model, zeg $\Struct{S}$. Volgens \olref{prop:equiv} geldt dan
  $\Struct{R} \equiv \Struct{S}$. Aangezien $\Struct{R}$ en $\Struct{S}$
  bovendien niet isomorf zijn, volgt dat de omkering van
  \olref[iso]{thm:isom} in het algemeen niet geldt.
\end{rem}

\end{document}
